{\bf \large The material of this chapter is specific to the Windows 64-bit version of STEPSS and its Java-based Graphic User Interface (GUI).

Before going any further, you have to read and accept the legal terms detailed at the beginning of this document.}

STEPSS is not restricted to Windows. Installers are published for Windows, macOS and Linux, RAMSES also builds on Linux and macOS with the GNU Fortran compiler, and stepss, the Python interface, is installed with {\tt pip} on both Windows and Linux. Those routes are documented online at\\
\hspace*{5ex}{\tt https://stepss.sps-lab.org}

Every version of STEPSS is published on its releases page, which is the place to download it from:\\
\hspace*{5ex}{\tt https://github.com/SPS-L/stepss-java-ui/releases}

Copies of the third-party installers mentioned in this chapter are also kept on a \hypertarget{Google}{Google drive}, to save navigating vendor web pages:\\
\hspace*{-7ex}{\tt https://drive.google.com/drive/folders/1HcUSD-FOx6192HURJnxltYJuKuNA8Vwz?usp=drive\_link}

{\bf \large Note that those third-party files are not part of STEPSS}.

%------------------------------------------------------------------------------------------------------
\section{Installing STEPSS}

STEPSS comes two ways. The {\bf installer} carries its own Java, so nothing else has to be installed and any Java already on the machine is neither needed nor disturbed. The {\bf archive}, {\tt stepss.jar}, is a single file that leaves nothing behind, and needs a Java you have installed yourself. Either way, everything required to run the simulations and display the results, executables and libraries included, travels with STEPSS.

\subsection{Using the installer (recommended)}

From the releases page, download the file for your system:
\begin{itemize}
\item Windows: {\tt STEPSS-\emph{version}.msi}
\item macOS (Apple Silicon): {\tt STEPSS-\emph{version}.dmg}
\item Linux (Debian, Ubuntu): {\tt stepss\_\emph{version}\_amd64.deb}
\end{itemize}

Open it and follow the installer. STEPSS is then started from the Start menu, from Applications, or from your desktop's list of applications, like any other program.

To remove it, uninstall it the way you would any other application: {\em Add or Remove Programs} on Windows, dragging it to the Trash on macOS, {\tt sudo apt remove stepss} on Debian and Ubuntu.

\subsection{Using the archive}

The archive route needs a 64-bit Java Runtime Environment (JRE), version~11 or later, installed on your computer. Any recent JRE or JDK is fine. To check whether you have Java, and which version:
\begin{enumerate}
\item open a Windows Command prompt
\item in the latter, enter the command:\\
\hspace*{5ex}{\tt java -version}

It should display a message similar to the one shown in Fig.~\ref{fig_install_java}.
\item close the command prompt.
\end{enumerate}

\begin{figure}[h!]
\centering
\includegraphics[scale=0.8]{java_check.jpg}
\caption{Checking the presence of the Java environment}\label{fig_install_java}
\end{figure}

If the command is not recognised, no Java is installed. In that case, install a 64-bit JRE or JDK as follows:
\begin{enumerate}
\item download a Java SE Development Kit (64-bit) from the \hyperlink{Google}{Google drive}, or a current release from {\tt https://adoptium.net/} or {\tt https://www.oracle.com/java/technologies/downloads/}
\item install by double-clicking on that file and following the instructions
\item for an easy execution, associate {\tt stepss.jar} with Java following the standard Windows procedure:
\begin{enumerate}
\item right-click on the {\tt stepss.jar} icon
\item select open with
\item choose always use this application.
\end{enumerate}
\end{enumerate}

There is no need to uninstall any Java version you already have, as long as it is 64-bit Java~11 or later.

Apart from the Java machine, this route does not require any installation in the usual meaning of the term.
Just download the {\tt stepss.jar} file from the releases page and drop it in any convenient place of the computer. The Windows desktop is a very convenient location.

Leave the archive intact; in particular do not uncompress it !

STEPSS is launched by merely double-clicking on {\tt stepss.jar} (or a shortcut to the latter).

When STEPSS is launched it creates a temporary working folder and copies all needed executables and libraries into that folder.

To remove STEPSS from your computer, just delete the archive file. STEPSS does not leave anything in your computer !

%--------------------------------------------------------------------------------------------------------
\section{Installing a Fortran Compiler}

{\bf \large This step can be skipped if you do not plan to develop new models for RAMSES}. Everything else in STEPSS works without it.

Compiling a model of your own produces a simulator built from the URAMSES kit that STEPSS carries, and that kit is a {\bf GNU Fortran} build. You therefore need {\tt gfortran}, GNU {\tt make} and OpenBLAS. The Intel oneAPI compiler and the Microsoft Visual Studio environment it relies on are no longer used and no longer need to be installed.

On Windows, these come from MSYS2:
\begin{enumerate}
\item install MSYS2 from {\tt https://www.msys2.org/}
\item open an MSYS2 terminal and install the toolchain:\\
\hspace*{5ex}{\tt pacman -S mingw-w64-x86\_64-gcc-fortran mingw-w64-x86\_64-openblas make}
\end{enumerate}

STEPSS looks for MSYS2 in {\tt C:$\backslash$msys64}, or wherever the {\tt MSYS2\_ROOT} environment variable points, because the MSYS2 installer does not add its compiler directory to the system PATH.

On Linux, install the same three with your package manager; on Debian and Ubuntu that is\\
\hspace*{5ex}{\tt sudo apt install gfortran make libopenblas-dev}\\
On macOS, use Homebrew:\\
\hspace*{5ex}{\tt brew install gcc openblas}

The kits STEPSS carries are specific to a gfortran ABI, and each platform's default compiler matches the kit shipped for it. If yours does not, STEPSS reports the exact compiler version to install rather than failing part way through a build.
